\documentclass[11pt,a4paper]{article}

\usepackage[utf8]{inputenc}
\usepackage[T1]{fontenc}
\usepackage[margin=1in]{geometry}
\usepackage{amsmath}
\usepackage{amssymb}
\usepackage{booktabs}
\usepackage{longtable}
\usepackage{array}
\usepackage{xcolor}
\usepackage{hyperref}

\hypersetup{
  colorlinks=true,
  linkcolor=black,
  citecolor=black,
  urlcolor=blue
}

\setlength{\parindent}{0pt}
\setlength{\parskip}{0.72em}
\renewcommand{\arraystretch}{1.16}

\newcommand{\code}[1]{\texttt{#1}}
\newcommand{\claimbox}[1]{%
  \vspace{0.5em}
  \noindent\fbox{%
    \begin{minipage}{0.94\linewidth}
    #1
    \end{minipage}}
  \vspace{0.5em}
}

\title{\textbf{DAD FieldWorks}\\
{\large A Reusable Evidence Contract Architecture for Validation-Aware Computational Electromagnetics Software}}
\author{Harun Aktas}
\date{Technical Whitepaper / Public Specification Draft\\Version 0.6 public\\May 2026}

\begin{document}

\maketitle

\begin{center}
\textbf{Claim status:} No external validation claim. No production readiness claim.\\
Copyright \copyright{} 2026 Harun Aktas. All rights reserved.\\
No license is granted for reuse, redistribution, or derivative work without written permission from the copyright holder.
\end{center}

\newpage

\begin{abstract}
DAD FieldWorks is a private research and development project being developed toward validation-aware RF and quantum hardware design workflows. This public specification draft describes a reusable evidence contract architecture for computational electromagnetics software: every computed result should carry metadata, evidence rows, gate rows, comparison rows, risk rows, required future work, and explicit claim boundaries.

Version 0.6 updates the public positioning toward RF and quantum-hardware-oriented microwave workflows, summarizes current public-safe project state from the private dashboard, and records the current PEC cavity bridge status conservatively. The project remains research and development. The public material is not external validation evidence, not a production readiness claim, and not a commercial solver equivalence claim.
\end{abstract}

\tableofcontents
\newpage

\section{Introduction}
\label{sec:introduction}

Computational electromagnetics work often produces useful numerical output before that output is ready to support a public claim. A frequency estimate, residual, field plot, time trace, or diagnostic comparison can be useful for engineering research while remaining insufficient for external validation, production use, or tool equivalence.

DAD FieldWorks treats that separation as an architectural problem. A result should not carry only numbers. It should carry the evidence context needed to understand what was computed, what was checked, what is still blocked, and what cannot yet be claimed.

\claimbox{\textbf{Public boundary.} This document describes public-safe architecture and status. It does not publish private solver source code, internal test suites, notebooks, private artifact files, private paths, or external organization information.}

\section{Project Direction}
\label{sec:project-direction}

DAD FieldWorks is being developed toward validation-aware RF and quantum hardware design workflows. The project explores electromagnetic modeling, RF PCB workflows, microwave structures, resonator and cavity analysis, solver research, and reproducible evidence records.

The current public quantum-hardware orientation is expressed through classical RF and microwave structures: feed lines, cavities, resonators, coupling regions, field build-up, ringdown, and electromagnetic evidence records. This is a direction of workflow relevance, not a claim that the project is a finished quantum hardware design tool.

The project does not currently claim qubit simulation, quantum processor modeling, Josephson junction modeling, Hamiltonian extraction, coherence-time prediction, GUI availability, external validation, production readiness, or commercial solver equivalence.

\section{Motivation}
\label{sec:motivation}

Electromagnetic solver research includes many layers that should remain visible:

\begin{itemize}
  \item solver family and numerical method,
  \item problem geometry, materials, and boundary conditions,
  \item grid, mesh, component, basis, and field-staggering conventions,
  \item source and observable definitions,
  \item analytical or alternative computational comparison rows,
  \item residual, convergence, and diagnostic metrics,
  \item ambiguity, risk, and required future work,
  \item decision classes, failure classes, and claim boundaries.
\end{itemize}

If these layers are compressed into a single pass/fail value, the result can be overstated. The evidence contract architecture keeps those layers explicit.

\section{Evidence Contract Architecture}
\label{sec:evidence-contract-architecture}

An evidence contract is a structured result model. It can be used for scalar eigenmode diagnostics, FDTD time-domain runs, FEM Helmholtz checks, Method of Moments workflows, RF network utilities, or comparison workflows.

\begin{longtable}{p{0.32\linewidth}p{0.58\linewidth}}
\toprule
\textbf{Record type} & \textbf{Purpose}\\
\midrule
\endhead
\code{SolverMetadata} & Solver family, method, case identity, version context, and scope.\\
\code{EvidenceRow} & Structured record for a value, metric, diagnostic fact, or provenance item.\\
\code{GateRow} & Named check with pass, fail, blocked, or not-ready state.\\
\code{ReferenceComparisonRow} & Analytical, manufactured, regression, or alternative computational comparison context.\\
\code{RiskRow} & Technical risk, ambiguity, mitigation, and residual uncertainty.\\
\code{RequiredFutureWorkRow} & Work required before stronger claims or next execution stages can open.\\
\code{ClaimBoundary} & Explicit statement of allowed and blocked claims.\\
\code{DecisionClass} & Controlled decision label for the current result.\\
\code{FailureClass} & Controlled failure or blocked-state label.\\
\code{ResultContractSummary} & Compact summary of decision, failure class, and claim flags.\\
\code{SolverEvidenceResult} & Top-level container for metadata, values, rows, comparisons, risks, future work, and boundaries.\\
\bottomrule
\end{longtable}

The architecture uses conservative defaults. Numerical success does not automatically change claim status.

\begin{center}
\begin{tabular}{ll}
\toprule
\textbf{Flag} & \textbf{Conservative default}\\
\midrule
\code{ExternallyValidatedQ} & \code{False}\\
\code{ProductionAllowedQ} & \code{False}\\
\bottomrule
\end{tabular}
\end{center}

\section{Current Public Technical Status}
\label{sec:current-public-technical-status}

Current public-safe project status, paraphrased from the private dashboard, is:

\begin{itemize}
  \item Active work is focused on records-only PEC cavity residual metric, generalized mass, and divergence or gauge bridge closure-readiness planning.
  \item The latest bridge construction records are internally consistent and claim safe at records level.
  \item Residual metric policy remains unresolved.
  \item Generalized mass policy remains unresolved.
  \item Divergence or gauge policy remains unresolved.
  \item Incidence compatibility closure and numerical consistency closure remain not ready.
  \item Native Yee eigenmode prototype planning remains not ready.
  \item Final vector, map, incidence, stencil, operator, matrix, residual, rank, nullspace, eigensolve, FDTD, CPML, validation, and production flags remain false for this bridge.
  \item C++ is infrastructure in progress and is not a production EM solver.
  \item No GUI exists.
\end{itemize}

The public site contains selected architecture notes, public status notes, generated diagnostic visuals, and public whitepaper material. It does not release private implementation details or internal artifacts.

\section{PEC Cavity Evidence Chain}
\label{sec:pec-cavity-evidence-chain}

The PEC cavity evidence chain is a central internal research thread because it exposes the difference between numerical output and claim readiness. It includes scalar Helmholtz diagnostics, FDTD component-family diagnostics, grid interpretation questions, boundary policy issues, residual policy issues, and native Yee planning questions.

\subsection{Reference Case}

The current public-safe reference case remains a rectangular PEC cavity:

\begin{center}
\begin{tabular}{ll}
\toprule
\textbf{Quantity} & \textbf{Value}\\
\midrule
\(a\) & \(0.08\ \mathrm{m}\)\\
\(b\) & \(0.084\ \mathrm{m}\)\\
\(d\) & \(0.084\ \mathrm{m}\)\\
\((m,n,p)\) & \((1,1,1)\)\\
\(c_0\) & \(299792458\ \mathrm{m/s}\)\\
Analytical reference frequency & \(3.143165987503105\ \mathrm{GHz}\)\\
\bottomrule
\end{tabular}
\end{center}

The analytical reference is scoped to this simple case. It does not create solver validation by itself.

\subsection{Scalar Helmholtz Refinement}

\begin{longtable}{ccccc}
\toprule
\textbf{Level} & \textbf{Grid} & \textbf{Frequency (GHz)} & \textbf{Relative error} & \textbf{Residual norm}\\
\midrule
\endhead
0 & \(\{9,8,8\}\) & 3.128307359058 & 0.0047272809 & \(2.63\times 10^{-15}\)\\
1 & \(\{18,16,16\}\) & 3.139012096545 & 0.0013215627 & \(1.29\times 10^{-14}\)\\
2 & \(\{27,24,24\}\) & 3.141247276062 & 0.0006104391 & \(1.68\times 10^{-14}\)\\
\bottomrule
\end{longtable}

These rows show bounded scalar diagnostic convergence. They are not a full vector Maxwell eigenmode claim.

\subsection{Same-Grid Scalar Diagnostic}

\begin{longtable}{ccccc}
\toprule
\textbf{Policy} & \textbf{Matrix size} & \textbf{Frequency (GHz)} & \textbf{Relative error} & \textbf{Residual norm}\\
\midrule
\endhead
1 & \(504\times 504\) & 2.778104015200 & 0.1161446687 & \(1.83\times 10^{-15}\)\\
2 & \(336\times 336\) & 3.122528486251 & 0.0065658325 & \(1.14\times 10^{-15}\)\\
3 & \(504\times 504\) & 3.126946915965 & 0.0051601066 & \(2.37\times 10^{-15}\)\\
\bottomrule
\end{longtable}

The same-grid scalar diagnostic remains inconclusive because mapping policy ambiguity is material. The numerical eigenvalue solves are clean, but interpretation is not closed.

\subsection{Internal Frequency Comparison}

\begin{longtable}{p{0.46\linewidth}cc}
\toprule
\textbf{Route or diagnostic} & \textbf{Frequency (GHz)} & \textbf{Relative error}\\
\midrule
\endhead
Analytical reference & 3.143165987503 & 0\\
Route A scalar Helmholtz Level 0 & 3.128307359058 & 0.004727281\\
Route A scalar Helmholtz Level 1 & 3.139012096545 & 0.001321563\\
Route A scalar Helmholtz Level 2 & 3.141247276062 & 0.000610439\\
Same-grid scalar Policy 2 & 3.122528486251 & 0.006565833\\
Same-grid scalar Policy 3 & 3.126946915965 & 0.005160107\\
Numerical dispersion diagnostic & 3.126477519038 & 0.005309445\\
FDTD EzHxFamily & 2.950641806786 & 0.061251675\\
FDTD ExFamily & 3.631559146814 & 0.155382554\\
\bottomrule
\end{longtable}

This is an internal evidence-path comparison. Component-resolved FDTD3D mode ambiguity remains unresolved. Internal candidate promotion remains blocked. Comparison execution remains blocked until ambiguity, residual policy, metadata policy, and acceptance gates close.

\section{RF And Quantum-Hardware-Oriented Direction}
\label{sec:rf-quantum-direction}

The quantum-hardware-oriented direction is currently classical electromagnetic and RF focused. Microwave resonators, feed lines, cavities, couplers, and field concentration regions are part of the classical microwave environment around many advanced devices.

At this stage, DAD FieldWorks public material supports discussion of:

\begin{itemize}
  \item microwave resonator diagnostics,
  \item cavity and resonator analysis,
  \item RF PCB workflow direction,
  \item field build-up and ringdown,
  \item evidence records and claim boundaries for electromagnetic results.
\end{itemize}

It does not claim qubit simulation, Josephson junction modeling, quantum state evolution, gate fidelity prediction, coherence-time prediction, or claim-bearing quantum processor assessment.

\section{FDTD Microwave Resonator Ringdown Diagnostic}
\label{sec:fdtd-resonator-ringdown}

The public FDTD microwave resonator ringdown visualization is a classical electromagnetic diagnostic. It shows a short pulse entering a feed line, coupling into a resonator-like region, building up an oscillating field pattern, and then decaying after the excitation has passed.

The left panel shows signed \(E_z\) for a 2D TMz FDTD diagnostic. Warm and cool colors represent opposite signs of the field component, not temperature. The upper right trace shows a local \(E_z\) probe response. The lower right trace summarizes field energy inside the resonator region.

The late-time decay is the ringdown. It is useful for communicating coupling, temporary energy storage, and resonator response in a classical microwave structure. It is not validation evidence and not a qubit simulation.

\section{C++ Evidence Contract Foundation}
\label{sec:cpp-foundation}

C++ is planned as a future production-core direction, but the current public-safe status is infrastructure only. The current C++ layer contains evidence-contract infrastructure and non-numerical demo scaffolding for result-contract shape. It is not a production EM solver.

The important public point is architectural: future solver wrappers should bind numerical results to structured metadata, evidence rows, gate rows, risk rows, required future work, and claim boundaries from the beginning.

\section{GUI Status}
\label{sec:gui-status}

No GUI exists. GUI planning should remain deferred until core workflows, supported operations, and claim boundaries are mature enough to expose safely.

\section{Claim Boundaries}
\label{sec:claim-boundaries}

This public draft does not claim:

\begin{itemize}
  \item external validation,
  \item production readiness,
  \item commercial solver equivalence,
  \item a complete quantum hardware solver,
  \item qubit simulation,
  \item Josephson junction modeling,
  \item quantum processor modeling,
  \item CPML support,
  \item open boundary support,
  \item GUI availability,
  \item a production C++ EM solver,
  \item a released private implementation.
\end{itemize}

\section{Roadmap}
\label{sec:roadmap}

\subsection{Near Term}

\begin{itemize}
  \item continue records-only PEC cavity residual metric, generalized mass, and divergence or gauge closure-readiness work,
  \item keep residual policy, mass policy, divergence or gauge policy, incidence closure, and numerical consistency closure explicit,
  \item preserve conservative validation and production flags,
  \item keep public RF and quantum-hardware-oriented wording tied to classical microwave diagnostics.
\end{itemize}

\subsection{Mid Term}

\begin{itemize}
  \item resolve residual and mass policy questions before stronger PEC cavity claims,
  \item plan native Yee curl-curl metadata and boundary requirements,
  \item harden RF PCB and resonator diagnostics with evidence rows,
  \item continue FDTD resonator and cavity evidence workflows without broadening claims,
  \item mature C++ evidence-contract wrappers without presenting them as solver releases.
\end{itemize}

\subsection{Long Term}

\begin{itemize}
  \item migrate stable research kernels into C++ only after evidence contracts and numerical gates are mature,
  \item plan GUI work only after supported operations and claim boundaries are stable,
  \item define external comparison pipelines only when traceable reference artifacts and acceptance policies exist,
  \item maintain claim-boundary review as a release discipline for every solver family.
\end{itemize}

\section{Conclusion}
\label{sec:conclusion}

DAD FieldWorks is organized around a simple discipline: computed results should carry evidence and claim boundaries. Version 0.6 updates the public whitepaper to reflect the current RF and quantum-hardware-oriented direction while keeping the project status conservative. The current work is meaningful research and architecture work, but major claims remain blocked until the corresponding evidence and gates exist.

\section{References}
\label{sec:references}

\begin{thebibliography}{9}

\bibitem{yee1966}
K. S. Yee,
``Numerical solution of initial boundary value problems involving Maxwell's equations in isotropic media,''
\emph{IEEE Transactions on Antennas and Propagation}, vol. 14, no. 3, pp. 302--307, 1966.

\bibitem{taflove}
A. Taflove and S. C. Hagness,
\emph{Computational Electrodynamics: The Finite-Difference Time-Domain Method}.
Artech House, third edition, 2005.

\bibitem{oberkampf}
W. L. Oberkampf and C. J. Roy,
\emph{Verification and Validation in Scientific Computing}.
Cambridge University Press, 2010.

\bibitem{roache}
P. J. Roache,
\emph{Verification and Validation in Computational Science and Engineering}.
Hermosa Publishers, 1998.

\end{thebibliography}

\end{document}
